\section{探索アルゴリズム}\label{sec:exploration}

探索フェーズは実行の実験的中核である:選ばれたアイディアと、それが
保持するメトリック契約 — 各々がその証拠となる正確な測定名を名指す、
反証可能な主張の実行レベルの宣言(\cref{sec:paper:contract}) — が与
えられたとき、探索は有界なノード予算を、それらの主張を評価可能にす
る測定へと変換せねばならない。ARI はこのフェーズを研究ステップ上の
最良優先木探索(best-first tree search, BFTS)として組織する。骨格は
AIDE に由来する。AIDE は LLM 駆動のエンジニアリングを、ハードコード
された draft/debug/improve 方針と単一スカラー目的の下での候補解の木
上の探索として枠づけた \cite{jiang2025aide}。ARI はその方針を、決定
的なフォールバックに支えられた LLM 主導の選択
(\cref{sec:exploration:node-selection})で置き換え、契約由来のカバレ
ッジと系譜ヒントで拡張を操舵し
(\cref{sec:exploration:expansion,sec:exploration:lineage-chaining})
— これは貪欲な単一目的方針について報告された局所最適のプラトー
\cite{jiang2025aide} への構造的な答えである — 各ノードを、その義務
とフィードバックが同じ契約から来るガード付き Reason-and-Act(ReAct)
エージェントとして実行する(\cref{sec:exploration:react})。締めくくり
の各小節は、探索が走る基盤を扱う:プラットフォームプローブ
(\cref{sec:exploration:platform-probe})、研究メモリ
(\cref{sec:exploration:memory})、そしてルートを播種する構想段階
(\cref{sec:exploration:ideation})である。

\subsection{形式化}\label{sec:exploration:formalization}

探索ループは木 $\Tree=(\Nodes, E)$ を維持し、各ノード
$n\in\Nodes$ は研究の一ステップ — アイディア改訂、コード変更、
ベンチマーク実行、分析など — を表す。ノードは以下を保持する:

\begin{itemize}
  \item 離散状態 $\statefn{n}\in\{\StateOpen,\StateEval,\StateDone,\StateFailed\}$、
  \item \texttt{NodeLabel} 列挙体から取られたカテゴリカルな
        \term{ラベル}、
  \item 軸ごとの評価指標を保持する自由形式 dict
        \texttt{node.metrics}、区別されたスカラー
        $\texttt{\_scientific\_score}\in[0,1]$ を含む、
  \item ノードが provisional なテキストではなく実測値を伴うかを
        示す Boolean \texttt{has\_real\_data}、ならびに
  \item 選択プロンプトおよび刈り込み判定で使われる帳簿フィールド
        \texttt{depth}。ARI は失敗ノードを再試行しない設計のため、
        \texttt{retry\_count} シグナルは一切提示されない。
\end{itemize}

\begin{figure}[t]
  \centering
  \resizebox{\columnwidth}{!}{% F02: BFTS exploration tree — snapshot of a small search tree, drawn in the
%      ARI house design language (_style.tex). One column per search depth;
%      every box is one tree node with its scalar score r. Peach boxes are
%      experimentation nodes; the pale-green box is the best-scoring node,
%      and the vermilion edges trace its lineage back to the root (the chain
%      the paper phase later consumes). Dashed gray edges mark the low-score
%      branch the selection policy stops picking. Node ids and scores come
%      from the frozen synthetic example recorded in data/F02_tree.meta.yaml
%      (illustrative, not a measured run).
% style.tikzstyles is loaded by shared/preamble.tex / the standalone wrapper;
% the ari* styles come from _style.tex, resolved below from whichever cwd
% this file is \input from (en|ja|zh builds, `make figure`, or the tikz dir).
\InputIfFileExists{../shared/figures/tikz/_style}{}{%
  \InputIfFileExists{shared/figures/tikz/_style}{}{%
    \InputIfFileExists{_style}{}{}}}%
\begin{tikzpicture}[aricanvas]
  % ---- depth-0 column: the root node --------------------------------
  \node[aribox exp]                    (n0) {\figlabel{F02.n0}};
  \node[ariphase, above=0.35cm of n0]  (h0) {\figlabel{F02.head.d0}};
  % ---- depth-1 column ------------------------------------------------
  \node[aribox exp, right=of n0]       (n1) {\figlabel{F02.n1}};
  \node[ariphase, above=0.35cm of n1]  (h1) {\figlabel{F02.head.d1}};
  % ---- depth-2 column ------------------------------------------------
  \node[aribox exp, right=of n1]       (n3) {\figlabel{F02.n3}};
  \node[ariphase, above=0.35cm of n3]  (h2) {\figlabel{F02.head.d2}};
  \node[aribox exp, below=of n3]       (n4) {\figlabel{F02.n4}};
  \node[aribox exp, below=0.9cm of n4] (n5) {\figlabel{F02.n5}};
  % ---- depth-1, low-score subtree -------------------------------------
  \node[aribox exp, left=of n5]        (n2) {\figlabel{F02.n2}};
  % ---- depth-3 column: the best-scoring node ---------------------------
  \node[aribox gate, right=of n4]      (n6) {\figlabel{F02.n6}};
  \node[ariphase] (h3) at (n6 |- h2)        {\figlabel{F02.head.d3}};
  % ---- expansion edges: lineage of the best node in accent -----------
  \draw[ariarrow accent] (n0) -- (n1);
  \draw[arielbow vh, draw=accent] (n1.south) to (n4.west);
  \draw[ariarrow accent] (n4) -- (n6);
  % ---- remaining expansion edges --------------------------------------
  \draw[ariarrow] (n1) -- (n3);
  \draw[arielbow vh, draw=arimuted, dashed] (n0.south) to (n2.west);
  \draw[ariarrow muted] (n2) -- (n5);
  % ---- sparse annotations (lifecycle reading) -------------------------
  \node[arinote, below=0.25cm of n6]   {\figlabel{F02.note.best}};
  \node[arinote, below=0.25cm of n5] (nlow) {\figlabel{F02.note.low}};
  % figure-level marker, bottom-left corner (away from any single node)
  \node[arinote] at (n0 |- nlow)       {\figlabel{F02.note.synthetic}};
\end{tikzpicture}
}
  \caption{小さな BFTS 探索木(合成・例示用)。列は探索深さ;各ボ
           ックスはスカラー報酬 $\nodescore{\cdot}$ を持つ 1 ノード
           である。強調された辺は最良ノードの系譜をルートまで辿る
           — 論文フェーズが消費する連鎖(\cref{sec:paper:evidence})
           である;破線の辺は選択方針が以降選ばなくなる低スコア分岐
           を示し、その履歴は木上に残る。}
  \label{fig:tree}
\end{figure}

\Cref{fig:tree} はこの種の小さな木を示す。軸レベル信号から
scientific score への集約は、\texttt{Evaluator} プロトコルの任意の
実装に委譲される;ARI は固定の線形結合を pin しない。プロトコルは
BFTS と下流のルブリック評価との唯一の継ぎ目であり、軸をスカラーへ
縮約する方法について検索エンジンを agnostic に保つ。

\paragraph{Run-loop 状態.}
外側ループの各反復は次のタプルを変換する:
\[
  \state \;=\; (\Frontier,\;\Pending,\;\expandcount,\;\Hist),
\]
ここで $\Frontier\subseteq\Nodes$ は拡張可能な完了ノード集合
(\StateDone{} / \StateFailed{}), $\Pending\subseteq\Nodes$ は実
行待ち(\StateOpen{})ノード集合, $\expandcount:\Nodes\to\Naturals$
は各フロンティアノードが何回展開されたかをカウント, $\Hist=(l_1,\ldots,l_t)$
は実行済みラベルのスライディングウィンドウ(長さ上限 $\histlen=20$、
\cref{sec:exploration:diversity}参照)である。初期状態は
$\state_0=(\emptyset,\{n_\mathrm{root}\},\mathbf{0},())$。

\paragraph{枝刈り述語.}
ハードな探索予算は次の述語で表される:
\begin{multline}\label{eq:prune}
  \prunepred(n;\,|\Nodes|) \;=\;
    \indicator{|\Nodes| \geq B_N}
    \;\lor\;
    \indicator{\depth{n} \geq B_D}
    \\\;\lor\;
    \sterilepred(n),
\end{multline}
ここで $B_N = \texttt{config.max\_total\_nodes}$, $B_D =
\texttt{config.max\_depth}$, sterile 述語
\begin{equation}\label{eq:sterile}
  \sterilepred(n) \;=\; \indicator{\,|\Delta_n| = 0\,}
\end{equation}
($\Delta_n$ は $n$ が追加・変更・削除したファイルの集合)
は親に対して 1 ファイルも新規アーティファクトを書き出さずに終了
したノードで真となる。run-loop がこのファイル差分を計
算し、結果を Boolean \texttt{\_sterile} フラグとしてノードに記録
する;\texttt{should\_prune} はこのフラグと 2 つの上限を読むだけ
である。呼び出し側が毎反復 live な $|\Nodes|$ を渡すことで,ノー
ド数の単一の真理源が保たれる。ステップ・コスト予算は
BFTS エンジン内ではなく呼び出し側のコストトラッカーによって別途
強制される。

\paragraph{Retire 述語.}
$\prunepred$ に加え、run-loop はより柔らかい \emph{フロンティア
retire 述語} $\retirepred$ を適用し、最も生産的なリードでなく
なった親 $f$ を $\Frontier$ から除去する:
\begin{equation}\label{eq:retire}
  \retirepred(f, c) \;=\;
    \underbrace{\indicator{\nodescore{c} > \nodescore{f}}}_{\text{規則 A}}
    \;\lor\;
    \underbrace{\indicator{\expandcount(f) \geq \maxexp}}_{\text{規則 B}},
\end{equation}
ただし $\maxexp = \texttt{config.max\_expansions\_per\_node}$
(既定値 $4$)。規則 A は子 $c$ が親 $f$ を上回った瞬間に $f$ を
retire し、規則 B は親ごとの予算を制限して探索を広げる。$\retirepred$ はノードを\emph{削除しない} — その履
歴は $\Tree$ に残る — 再展開対象から外すだけである。

外側ループ 1 反復の段階分解(内部展開サブループ、並列 ReAct バ
ッチ、実行後 retire)は \cref{fig:bfts-state} に示す。明示的な
擬似コードは \cref{sec:exploration:run-loop} の
\cref{alg:bfts-run-loop} を参照。

\begin{figure}[t]
  \centering
  \resizebox{0.85\columnwidth}{!}{% F15: BFTS state-machine view of one outer run-loop iteration, set in the
% ARI house design language (_style.tex). Single top-down spine; the two
% iteration scopes are dashed ariloop enclosures — the inner expansion
% sub-loop (steps 1-3) and the outer parallel-run loop (steps 4-7) —
% joined through the white state pools they communicate over (frontier
% F at the top, pending P between them). Pale-green boxes are the
% predicate gates (Pi, sigma, rho); the vermilion edges (the single
% accent role of this figure) are the side effects fired when a gate
% predicate evaluates true; the dotted gray margin path is the reflux
% into the next outer iteration. Auxiliary state (e, H) stays in the
% caption, exactly as before; per-edge state-bump annotations were
% folded into the caption to keep annotation density low.
%
% style.tikzstyles is loaded by the preamble (inlined) and by the
% standalone wrapper; the refined ari* styles used below come from
% _style.tex in this directory (gate T2: all widths live in the style
% files, never inline).
%
% The build cwd differs per entry point, so resolve _style.tex from each
% known prefix: report/<lang>/ (the language PDFs \input figures as
% ../shared/figures/tikz/<fig>.tex), report/ (make figure / render_tikz),
% the repo root, then this directory itself (direct cwd or TEXINPUTS).
\InputIfFileExists{../shared/figures/tikz/_style.tex}{}{%
  \InputIfFileExists{shared/figures/tikz/_style.tex}{}{%
    \InputIfFileExists{report/shared/figures/tikz/_style.tex}{}{%
      \InputIfFileExists{_style.tex}{}{%
        \errmessage{F15: cannot locate figures/tikz/_style.tex}}}}}
\begin{tikzpicture}[aricanvas]
  %% Heading: the state tuple the run loop mutates.
  \node[ariphase] (hd) {\figlabel{F15.band.state}};
  \pic at ([xshift=-2.5mm]hd.west) {ariglyph flask};

  %% Spine. White boxes are shared state pools, peach boxes the
  %% experimentation-phase steps, pale-green boxes the predicate gates.
  \node[aribox, below=0.4cm of hd]                    (frontier) {\figlabel{F15.frontier}};
  \node[ariboxwide, fill=phaseexp, below=0.9cm of frontier] (pick) {\figlabel{F15.pick_best}};
  \node[aribox gate, below=of pick]                   (prune)    {\figlabel{F15.prune}};
  \node[aribox exp, below=of prune]                   (expand)   {\figlabel{F15.expand}};
  \node[aribox, below=0.8cm of expand]                (pending)  {\figlabel{F15.pending}};
  \node[ariboxwide, fill=phaseexp, below=0.9cm of pending]  (batch) {\figlabel{F15.select_run}};
  \node[ariboxwide, fill=phaseexp, below=of batch]    (react)    {\figlabel{F15.run_parallel}};
  \node[aribox gate, below=of react]                  (sterile)  {\figlabel{F15.sterile}};
  \node[aribox, below=of sterile]                     (append)   {\figlabel{F15.append_frontier}};
  \node[aribox gate, below=of append]                 (retire)   {\figlabel{F15.retire}};

  %% Gate side effects (fire when the predicate is true).
  \node[aribox, right=of prune]   (dropf) {\figlabel{F15.prune_drop}};
  \node[aribox, right=of sterile] (clamp) {\figlabel{F15.clamp}};
  \node[aribox, right=of retire]  (dropr) {\figlabel{F15.retire_drop}};

  %% Loop enclosures: inner expansion sub-loop / outer parallel-run loop.
  \node[ariloop, fit=(pick)(prune)(expand)]                  (loopin)  {};
  \node[ariloop, fit=(batch)(react)(sterile)(append)(retire)] (loopout) {};
  \node[arinote, anchor=south west] at ([xshift=6mm]loopin.north)  (ttlin)  {\figlabel{F15.lane.inner}};
  \node[arinote, anchor=south west] at ([xshift=6mm]loopout.north) (ttlout) {\figlabel{F15.lane.outer}};
  \pic at ([xshift=-2.5mm]ttlin.west)  {ariglyph loop};
  \pic at ([xshift=-2.5mm]ttlout.west) {ariglyph loop};

  %% Step badges 1..7 (macro flow order; replaces edge labels).
  \node[aribadge] at (pick.north west)    {1};
  \node[aribadge] at (prune.north west)   {2};
  \node[aribadge] at (expand.north west)  {3};
  \node[aribadge] at (batch.north west)   {4};
  \node[aribadge] at (react.north west)   {5};
  \node[aribadge] at (sterile.north west) {6};
  \node[aribadge] at (retire.north west)  {7};

  %% Main top-down flow.
  \draw[ariarrow] (frontier) -- (pick);
  \draw[ariarrow] (pick)     -- (prune);
  \draw[ariarrow] (prune)    -- (expand);
  \draw[ariarrow] (expand)   -- (pending);
  \draw[ariarrow] (pending)  -- (batch);
  \draw[ariarrow] (batch)    -- (react);
  \draw[ariarrow] (react)    -- (sterile);
  \draw[ariarrow] (sterile)  -- (append);
  \draw[ariarrow] (append)   -- (retire);

  %% Gate-fired side effects: the one accent role in this figure.
  \draw[ariarrow accent] (prune)   -- (dropf);
  \draw[ariarrow accent] (sterile) -- (clamp);
  \draw[ariarrow accent] (retire)  -- node[arinote, above] {\figlabel{F15.yes_ab}} (dropr);

  %% Reflux: retire's "no" arm starts the next outer iteration.
  \draw[arifeedback] (retire.west) -- ([xshift=-1.6cm]retire.west)
                     |- (frontier.west);
  \node[arinote] at ([xshift=-1.6cm]pending.west) {\figlabel{F15.no_next}};
\end{tikzpicture}
}
  \caption{BFTS 外側ループ 1 反復の状態機械ビュー(上から下への
           単一フローチャート)。ステップ 1--3 が内部展開サブルー
           プ、ステップ 4--7 が外側並列実行ループ。黄色のひし形は
           述語ガード($\prunepred,\sterilepred,\retirepred$)、
           右側の赤破線ボックスは各述語が起こす副作用(フロンティ
           ア除去 / スコアクランプ / retire)。図中に描かない補助
           状態:親ごとの展開カウンタ $\expandcount(\cdot)$ はス
           テップ 3 で増分されステップ 7 の $\retirepred$ で参照
           される。ラベル履歴 $\Hist$(窓 $\histlen$)はステップ
           5 の \texttt{record\_run} で追記され、ステップ 4 で多
           様性ボーナス $\divbonus{\cdot}$ の計算に用いられる。
           \cref{fig:bftsflow}(同じループを LLM プロンプト込みの
           細粒度制御フローとして示す)と比較されたい。}
  \label{fig:bfts-state}
\end{figure}

\subsection{Run-loop アルゴリズム}\label{sec:exploration:run-loop}

\Cref{alg:bfts-run-loop} は反復を明示的な擬似コードとして書き、
\cref{fig:bftsflow} は同じステップをプロンプトファイル込みの制御
フロー背骨として描く。内側 \textsc{while}
ループは空きワーカースロットを 1 展開ずつ埋め(ワーカーは次の展
開提案前に必ず走り始める)、外側ブロックは
\texttt{select\_next\_node} を 1 回 1 ノードずつ繰り返し呼んで
ワーカー上限までバッチを組み立て、そのバッチを並列実行する。実
行後ブロックは (a) 実行されたラベルを $\Hist$ に追記して多様性
ボーナスを実態に整合させ、(b) $\retirepred$ の規則 A / B を適用
する。

\begin{algorithm}[t]
\small
\caption{BFTS run-loop(外側ループ 1 反復)。}
\label{alg:bfts-run-loop}
\begin{algorithmic}[1]
\Require 状態 $\state=(\Frontier,\Pending,\expandcount,\Hist)$;
         設定 $(B_N, B_D, \maxexp)$ とワーカー上限
         $\texttt{max\_parallel}=\min(\texttt{max\_parallel\_nodes},4)$
\Ensure  更新後の状態 $\state'$
\State \Comment{--- 内部展開サブループ ---}
\While{$\Frontier \neq \emptyset \land |\Pending| < \texttt{max\_parallel} \land |\Nodes| < B_N$}
  \State $f \gets \selectLLM(\Frontier)$ \Comment{\texttt{select\_best\_to\_expand}}
  \If{$\prunepred(f;\,|\Nodes|)$}
    \State $\Frontier \gets \Frontier \setminus \{f\}$;\;\;\textbf{continue}
  \EndIf
  \State $c \gets \texttt{expand}(f,\,\texttt{depth\_note},\,\texttt{budget\_note})$
  \State $\Pending \gets \Pending \cup \{c\}$;\;\;$\expandcount(f) \mathrel{+}{=} 1$
\EndWhile
\State \Comment{--- 外側並列実行バッチ ---}
\State $B \gets \emptyset$
\While{$|B| < \min(\texttt{max\_parallel},\,|\Pending|)$}
  \State $n \gets \selectLLM(\Pending)$ \Comment{\texttt{select\_next\_node}: 1 ノード}
  \State $\Pending \gets \Pending \setminus \{n\}$;\;\;$B \gets B \cup \{n\}$
\EndWhile
\ForAll{$n \in B$ \textbf{並列}}
  \State $n' \gets \mathrm{ReAct}(n)$ \Comment{\texttt{AgentLoop.run}、失敗あり}
  \State $\Hist \gets (\Hist \mathbin{\Vert} \mathrm{label}(n'))[\,-\histlen{:}\,]$ \Comment{\texttt{record\_run}}
  \State $\Frontier \gets \Frontier \cup \{n'\}$
  \ForAll{$p \in \Frontier : p = \parentof{n'}$} \Comment{規則 A}
    \If{$\nodescore{n'} > \nodescore{p}$}
      \State $\Frontier \gets \Frontier \setminus \{p\}$
    \EndIf
  \EndFor
  \ForAll{$p \in \Frontier$} \Comment{規則 B}
    \If{$\expandcount(p) \geq \maxexp$}
      \State $\Frontier \gets \Frontier \setminus \{p\}$
    \EndIf
  \EndFor
\EndFor
\State \Return $\state' = (\Frontier,\Pending,\expandcount,\Hist)$
\end{algorithmic}
\end{algorithm}

\subsection{ノード選択(LLM 主導)}\label{sec:exploration:node-selection}

\begin{figure}[t]
  \centering
  \resizebox{\columnwidth}{!}{% F03: signals consumed by the LLM-driven node selector.
% Replaces the earlier closed-form utility figure: ARI does not pin
% a scalar utility; instead a structured candidate description is
% handed to the LLM and an index is parsed back. See F11 for the full
% control flow; this figure focuses on the *information* that flows in.
% style.tikzstyles is loaded by shared/preamble.tex / standalone wrapper.
\begin{tikzpicture}[node distance=0.5cm and 1.6cm]

  % per-node signals (left column), all stations for paper-style typography
  \node[station]                                       (real)   {\figlabel{F3.has_real}};
  \node[station, below=of real]                        (metrics){\figlabel{F3.metrics}};
  \node[station, below=of metrics]                     (depth)  {\figlabel{F3.depth_signal}};
  \node[station, below=of depth]                       (label_st){\figlabel{F3.retry}};
  \node[station, below=of label_st]                    (summ)   {\figlabel{F3.summary}};
  \node[station, fill=ari-pink!22, below=of summ]      (div)    {\figlabel{F3.div}};

  % LLM selector (centre) — wider station emph
  \node[station emph big, right=2.2cm of metrics, yshift=-1.5cm]
                                                       (sel)    {\figlabel{F3.selector}};
  % the consumed prompt file (above the selector)
  \node[station, fill=ari-yellow!30, draw=ari-orange,
        above=0.6cm of sel]                            (prom)   {\figlabel{F3.prompt}};

  % output (right) — single index back to caller
  \node[station emph big, right=2.2cm of sel]          (idx)    {\figlabel{F3.index}};

  % deterministic fallback when LLM index is unparseable
  \node[station, dashed, draw=ari-vermilion, below=1.0cm of sel]
                                                       (fb)     {\figlabel{F3.fallback}};

  % --- numbered step badges --------------------------------------
  \node[numbered step, above left=0.05cm and -0.2cm of real.north west] {1};
  \node[numbered step, above left=0.05cm and -0.2cm of sel.north west]  {2};
  \node[numbered step, above left=0.05cm and -0.2cm of idx.north west]  {3};

  % --- edges -----------------------------------------------------
  \foreach \n in {real,metrics,depth,label_st,summ}
    \draw[arrow] (\n.east) -- (sel.west);
  \draw[arrow dashed] (div.east) to[bend right=8] (sel.west);
  \draw[arrow dashed] (prom.south) -- (sel.north);
  \draw[big arrow] (sel.east) -- (idx.west)
       node[edge label, font=\sffamily\footnotesize]{\figlabel{F3.parse}};
  \draw[big arrow, draw=ari-vermilion] (sel.south) -- (fb.north)
       node[edge label, font=\sffamily\footnotesize]{\figlabel{F3.unparseable}};
  \draw[big arrow, draw=ari-vermilion] (fb.east) to[bend right=22] (idx.south);

  % --- background grouping for input column ----------------------
  \begin{scope}[on background layer]
    \node[stage container, fit=(real)(metrics)(depth)(label_st)(summ)(div)] (col) {};
    \node[stage title, above=4pt of col.north west, anchor=south west] {\figlabel{F3.band.signals}};
  \end{scope}
\end{tikzpicture}
}
  \caption{次に操作するノードを選ぶ LLM に渡される選択シグナル。
           入力は構造化された候補記述に連結される;ソフトな
           \texttt{diversity\_bonus} はタイブレーカであり、主シ
           グナルではない。}
  \label{fig:nodescore}
\end{figure}

ARI は意図的にノードのランキングを閉形式のスカラー効用に\emph{
縮約しない} — これは ARI の BFTS を AIDE のハードコードされた貪欲
方針 \cite{jiang2025aide} から最も分かつ設計上の選択である。代わ
りに、\texttt{select\_next\_node}
(「次のエージェントステップが実行すべき pending ノードはどれ
か?」)と \texttt{select\_best\_to\_expand}(「拡張するに最も値
する完了ノードはどれか?」)の双方が候補記述リストを LLM に渡し、
0 始まりの単一インデックスを解析して受け取る — 各呼び出しはちょ
うど 1 ノードを選ぶ。各ノード $n$ の候補記述(\cref{fig:nodescore})
は次の形式の一行:
\begin{quote}\small\texttt{[i] id=..., depth=..., label=..., has\_real\_data=..., metrics=\{...\}, summary=...}\end{quote}
実行選択側(\texttt{select\_next\_node})は、現在ボーナス対象のノー
ドに対してこの行末へ \texttt{diversity\_bonus=+0.05} を付加する;
\texttt{select\_best\_to\_expand} は付加しない。これを消費するプ
ロンプトは \cref{prompt:bfts_select}(選択)および
\cref{prompt:bfts_expand_select}(拡張対象)に逐語掲載する。プロ
ンプトが列挙する選択基準は:
\texttt{has\_real\_data=True} で強い metrics を持つノードは高価
値;metrics は全体論的に多目的に重み付け;優れた結果を持つ深い
ノードは継続に値する;\texttt{diversity\_bonus} はソフトなタイ
ブレーカとしてのみ扱う、である。\texttt{label} フィールドは
\texttt{select\_best\_to\_expand} に示される情報と一致しており、
誤解を招く「low retry を優先」基準は存在しない。

\subsubsection*{多様性ボーナス}\label{sec:exploration:diversity}

多様性ボーナス \texttt{BFTS.diversity\_bonus} は最近のラベルを
\texttt{\_recent\_label\_history}(\texttt{record\_run} で埋まる
スライディング窓)で追跡する。当該窓内で最頻ラベルに対し相対的に
\emph{過少表現}のラベルを持つノードに対し、関数は次を返す:
\begin{equation}\label{eq:diversity}
  \mathrm{div}(n) =
    \begin{cases}
      +0.05 & \mathrm{cnt}(n) \le \tfrac{1}{2}\,\mathrm{cnt}_{\max},\\
      0     & \text{それ以外,}
    \end{cases}
\end{equation}
ただし $\mathrm{cnt}(n)$ は窓内における $\mathrm{label}(n)$ の出
現回数で,$\mathrm{cnt}_{\max}=\max_\ell\mathrm{cnt}(\ell)$ である。
$0.05$ という定数は意図的に小さい:scientific score が依然として
ランキングを支配し、ボーナスはほぼ同点のものを破るのみである。

\subsubsection*{LLM フォールバック}\label{sec:exploration:fallback}

LLM が解析不能なインデックスを返した場合、
\texttt{select\_next\_node} は \texttt{\_select\_fallback} ヘルパ
(候補を \texttt{\_fallback\_score} 式で順位付けする)による決定
的なランキングへフォールバックする。デフォルトの式は
\begin{equation}\label{eq:fallback}
  \fallbackscore{n} \;=\; \nodescore{n} + \mathrm{div}(n),
\end{equation}
ここで $\nodescore{n}\equiv\texttt{\_scientific\_score}(n)$ はノー
ドのスカラー報酬(\cref{sec:notation})である。
\texttt{\_select\_fallback} は \texttt{has\_real\_data=True} のノ
ードがあればそれを優先し、次いで $\fallbackscore{\cdot}$ 最大の候
補を返す。これにより、LLM 呼び出しが degrade してもループは必ず
前進する。

\subsubsection*{評価層の設定}\label{sec:exploration:config-layers}

\cref{sec:exploration:formalization,sec:exploration:node-selection,sec:exploration:fallback}
で登場する 4 つの評価面は独立に設定可能である(デフォルト値は上式
を再現するため、未変更の設定は no-op となる)。

\begin{description}
  \item[\textbf{合成式.}]
        \code{evaluator.composite} で指定。
        $\{a_k\}\!\to\!\texttt{\_scientific\_score}$ の縮約として、
        加重調和平均(デフォルト)、加重算術平均、ボトルネック型の
        \code{weighted\_min}、加重幾何平均から選択。調和・幾何
        平均は分母を有限に保つために $\epsilon=0.01$ で 0 軸を
        フロアする。
  \item[\textbf{フロンティアスコア戦略.}]
        \code{bfts.frontier\_score} で指定。
        \cref{eq:fallback} は \code{scientific\_plus\_diversity}
        に相当。他の選択肢は \code{scientific\_only}
        ($\mathrm{div}$ を落とす);\code{depth\_penalized} は
        $\mathrm{div}(n)$ に $-\lambda\,\depth{n}$ を加える
        ($\lambda=$\code{depth\_penalty\_lambda});
        \code{ucb\_like} は $\mathrm{div}(n)$ に
        $c_{\mathrm{ucb}}\sqrt{\log N / (\expandcount(n)+1)}$ を加える
        ($c_{\mathrm{ucb}}=$\code{ucb\_c}、
        $N = \sum_{n'}\expandcount(n')+|\Frontier|$)。
  \item[\textbf{軸セット.}]
        \code{evaluator.axis\_mode} で指定。
        \texttt{dynamic}(デフォルト)は汎用 5 軸フロア + 有効
        rubric + \texttt{idea.json} のプランキーワードから自動構築。
        \texttt{legacy} は 5 軸固定。\texttt{custom} は
        $\{(\mathrm{name},\mathrm{description},w)\}$ のリストを
        そのまま judge LLM に渡す。
  \item[\textbf{選択プロンプト.}]
        \code{bfts.select\_prompt} および
        \code{bfts.expand\_select\_prompt} で指定。
        \cref{alg:bfts-run-loop} の 2 つの $\selectLLM$
        呼び出しで使うテンプレートを、\code{FilesystemPromptLoader}
        キーで差し替え可能。selection テンプレートは
        \code{experiment\_goal}、\code{memory\_context}、
        \code{candidates} の placeholder を受け取り、expansion-target
        テンプレートは \code{experiment\_goal} と
        \code{candidates} を受け取る。
\end{description}

\subsection{拡張(1 呼び出しにつき 1 子)}\label{sec:exploration:expansion}

\texttt{expand} 関数は 1 回の呼び出しで\emph{高々 1 つ}の新
しい子を生成する(事前バッチ化なし)。同じ親に対して複数の子が欲しい
呼び出し側は
\texttt{expand} を反復的に呼び、これまでに生成された子を
\texttt{existing\_children} 引数経由で渡すことで、LLM が重複方向
を提案するのを避ける。

各新子のラベルはハードコードされたテンプレートではなく LLM の判
断で決まる;拡張プロンプト(\cref{prompt:bfts_expand} に逐語掲
載)には以下が渡される:

\begin{itemize}
  \item 親の状態(\texttt{succeeded} /
        \texttt{failed/no-real-data})と
        \texttt{\_scientific\_score};
  \item 親の構造化された \texttt{node\_report}
        (\code{delta\_vs\_parent} / \texttt{concerns} /
        \texttt{hints}、node-report ビルダが生成)が存在する場合、
        planner が特定の弱点を狙えるよう含める;
  \item 同じ深さの兄弟スコア、ならびにルートから親までの先祖ス
        コア;
  \item ツリー全体の多様性指標(これまでに見たユニークラベル、
        深さ分布);ならびに
  \item この親で既に生成された子のラベル(重複提案を避けるた
        め)。
\end{itemize}

子は \StateOpen{} で生成され、エージェントループに実行のため渡
され(\cref{sec:exploration:react})、すべての軸が埋まれば
\StateDone{} に、エージェントループが例外を投げれば
\StateFailed{} に遷移する。

\paragraph{カバレッジ操舵.}
拡張対象セレクタは本質的にブランチを跨ぐスケジューラ — 各ブラン
チのスコア、metrics、要約のすべてを見る — であるが、スコアだけで
は実行が宣言した主張について何も語らず、ある実際の実行では 10 ノ
ードの木が同じ目玉実験の 10 通りの変種を生む一方で、宣言された機
構の主張には専用の実験が一つも当てられなかった。そこで
\texttt{select\_best\_to\_expand} に渡される目標テキストは、メト
リック契約から導かれる実行レベルの \term{主張カバレッジ}ブロック
で拡張される:どの宣言済み主張が実行のどこかですでに裏づけ測定を
持ち、どれが未カバーのままかを示し、「未カバーの主張をカバーす
る」が\emph{どの}ノードを拡張するかを左右できるようにする。ある
主張が要求する測定名はその \term{契約証拠名}である。カバレッジは
二点で保守的に計算される:ある主張がカバー済みとされるのは、その
要求する名前の\emph{過半数}が出力されたときに限る — これは最終ゲ
ートの任意名一致規則よりも意図的に厳しい。操舵においては偶然に共
有された一つの汎用名が専用実験を永遠に抑止しかねない一方、未カバ
ーと誤る方は高々冗長な測定を冒すだけだからである — そして、評価器
が実データありと認めたノードのみが寄与し、ゲートと整合する。これ
により、欠陥のあるブランチの単なる測定\emph{名}が兄弟に「すでにカ
バー済み」と告げることはできない。同じ状態はノードごとに再計算さ
れ、実行中の各エージェントにも届く(\cref{sec:exploration:react})。
\emph{計算された}証拠を単一の系譜に沿って連鎖させる対のヒントは
\cref{sec:exploration:lineage-chaining} の主題である。

\subsection{枝刈りと停止}\label{sec:exploration:termination}

\texttt{should\_prune} は枝刈り述語 $\prunepred$(\cref{eq:prune})
を実装する:総ノード上限、深さ上限(\texttt{max\_depth} 設定フィー
ルド)、および sterile フラグである。
フロンティア retire 述語 $\retirepred$(\cref{eq:retire})
は各ノードの完了後に適用される:ノードを削除せず、その履歴は
$\Tree$ に残るが、フロンティアのプールが縮むため、同じ親を無限に
掘り続けるのではなく探索が広がる。

ループは次のいずれかが満たされれば停止する:

\begin{itemize}
  \item グローバル \texttt{max\_total\_nodes} 上限到達;
  \item フロンティアの\emph{すべての}ノードが $\Pi$ に該当(深さ
        上限、sterile、または予算枯渇)し pending プールも空である;
  \item コストトラッカーが強制する呼び出しごとのステップ / コスト
        予算の枯渇;
  \item \term{停滞検出器} \texttt{detect\_stagnation} が、
        \texttt{\_scientific\_score} の移動最大値がスライディング
        窓内で改善しないことを観測;
  \item 系譜判定モジュールからの \texttt{LineageDecision} が当該
        分岐を明示的に停止(\cref{prompt:lineage_decision} 参照)。
\end{itemize}

これらのうちどれが実際に最も頻出するかは経験的問いであり、凍結
チェックポイントが揃い次第 \cref{chapter:eval} で答える;本レポ
ートはその分布を主張しない。

\subsection{ReAct ドライバ}\label{sec:exploration:react}

\begin{figure}[t]
  \centering
  \resizebox{\columnwidth}{!}{% F04: ReAct loop — Think / Act / Observe / Reflect with numbered stages,
%       paper-style typography.
% style.tikzstyles is loaded by shared/preamble.tex / standalone wrapper.
\begin{tikzpicture}[node distance=2.0cm and 2.5cm]

  \node[state init]                            (think) {\figlabel{react.think}};
  \node[state, right=of think]                 (act)   {\figlabel{react.act}};
  \node[state, below=of act]                   (obs)   {\figlabel{react.observe}};
  \node[state final, below=of think]           (refl)  {\figlabel{react.reflect}};

  % numbered step badges, perched on each state's upper-left corner
  \node[numbered step, above left=0.05cm and -0.15cm of think.north west] {1};
  \node[numbered step, above left=0.05cm and -0.15cm of act.north west]   {2};
  \node[numbered step, above left=0.05cm and -0.15cm of obs.north west]   {3};
  \node[numbered step, above left=0.05cm and -0.15cm of refl.north west]  {4};

  % external annotations (yellow = memory, orange = phase-filtered tools)
  \node[station, fill=ari-orange!22, draw=ari-vermilion,
        above=0.6cm of act]                    (tools) {\figlabel{F4.toolset}};
  \node[station, fill=ari-yellow!30, draw=ari-orange,
        left=0.6cm of think]                   (mem)   {\figlabel{F4.memcontext}};

  % --- big arrows (paper-style) ---
  \draw[big arrow] (think) -- (act)   node[edge label, font=\sffamily\footnotesize]{\figlabel{F4.tool_call}};
  \draw[big arrow] (act)   -- (obs)   node[edge label, font=\sffamily\footnotesize]{\figlabel{F4.execute}};
  \draw[big arrow] (obs)   -- (refl)  node[edge label, font=\sffamily\footnotesize]{\figlabel{F4.parse}};
  \draw[big arrow] (refl)  -- (think) node[edge label, font=\sffamily\footnotesize]{\figlabel{F4.update}};
  \draw[arrow dashed] (obs) to[bend left=18] (think);

  \draw[arrow dashed] (tools.south) -- (act.north);
  \draw[arrow dashed] (mem.east)    -- (think.west);

  % budget annotation
  \node[stage title, below=0.3cm of refl] {\figlabel{F4.maxsteps}};
\end{tikzpicture}
}
  \caption{ReAct ループ。実線が主サイクル;破線のバックエッジは
           観測がすでに開いた問いを閉じた場合の早期離脱を表す。}
  \label{fig:react}
\end{figure}

各ノードの拡張は \texttt{AgentLoop} クラスの ReAct エージェント
ループ \cite{yao2023react} を走らせる(\cref{fig:react}):モデルは
蓄積された文脈について推論し、ツール呼び出しを発行し、構造化され
た結果を観測し、これを繰り返す。最大ステップ数は
\texttt{MAX\_REACT\_STEPS = 80} で、インスタンスごとに
\texttt{max\_react\_steps} キーワードで上書きできる。別ガード
\texttt{MIN\_TOOL\_CALLS = 2} が trivially 短い軌跡を防ぐ。エー
ジェントのシステムプロンプトは \cref{prompt:agent_system} に逐
語掲載する。

エージェントが利用できるツール集合はツールマネージャがフェーズ
別にフィルタする;例えば探索フェーズは論文執筆ツールをマスクし、
BFTS 拡張が草稿にショートカットされるのを防ぐ。一部の MCP ツー
ルは \texttt{ari-core} 自身用に予約されており(memory CoW 操
作)、LLM からは隠されるため、モデルが任意のノード ID を設定し
てメモリスキルのコピーオンライトチェックを迂回することはできな
い。

\subsubsection*{ノード開始時の義務}

ノードは空白のプロンプトから始まるわけではない。最初のモデル呼び
出しの前に、ループは決定的な文脈を注入する:\cref{sec:exploration:memory}
の \term{作業文脈} — 実験の固定核とその先祖の結論 — と、実行のメ
トリック契約が存在するようになれば、ドメイン中立な言葉で、最終的
な主張が単に尤もらしいのではなく科学的に妥当であるためにノードが
何をせねばならないかを述べる \term{契約義務}メッセージである:目
玉数値が用いる問題サイズで、独立した参照に対し正しさを検証し、残
差を測定として記録する;正規化の分母に使われる天井は決して\emph{
ハードコードせず} — \emph{測定する};各天井と正しさ検査の来歴をタ
グ付けし、それらが実行されたことをゲートが確認できるようにする;
目玉メトリックがその生のオペランドからどう再計算されるかを宣言す
る;そして各宣言済み主張を評価可能にする測定を、\emph{正確な}宣
言名のもとで出力する — 否定的な結果は許容されるが、証拠のない主張
は許容されない。このテキストはドメイン知識を一切持たない;エージ
ェントは自らのドメインが要請するとおりに各義務を果たす(スループ
ット天井ならマイクロベンチマーク、精度天井ならベースライン実行、
数値法なら解析的検査)。

実行レベルの付録 3 つが同じメッセージに同乗する:ノードの主張カバ
レッジ状態(\cref{sec:exploration:expansion})は、実現可能に測定で
きる未カバー主張を 1 つか 2 つ優先するよう求める;
\cref{sec:exploration:lineage-chaining} の継承データノート;および
\cref{sec:exploration:platform-probe} のプラットフォームノートで
ある。これらの実行レベル不変量は \emph{pin} される:スライディン
グ文脈窓はこれらを — survey、アイディア生成、metric-spec ツール結
果とともに — 常に保持し、古いターンは脱落する。ルートでは、文脈構
築時に契約がまだ存在しないため、義務は契約鋳造ツール呼び出しのハ
ンドラがループ途中で注入する。こうしてどのノードも宣言済み主張を
知らずに走ることはなく、義務を二度受け取るノードもない。

\subsubsection*{出力フィードバック}

測定は単一の構造化出力ツール(\texttt{emit\_results})を通じてノー
ドを離れる。これはノードの結果ファイルを書き、メトリック契約に対
し検証し、まだ欠けている必須証拠を名指す \term{契約警告}を返す。
初期の実行はまさにこの地点で失敗の形を露呈した:警告はツール結果
の内部に届き、ハーネスは最初の完了ジョブの直後にノードを強制終了
し、満たされない義務は行動するステップを残されなかった — 観測され
たあるノードは 80 ステップのうち 66 ステップを未使用のまま終わっ
た。そこでループは出力時点のフィードバックを実行可能なターンへ変
える:出力が契約警告を伴うとき、継続を促すナッジ — 欠けている測
定、残りステップ、そしてそれらを今すぐ実行して再出力せよ(再出力
は上書きする)、あるいはこのノードでは実現不能な場合のみ結論せよ、
という指示 — を注入し、ノードを次に述べる強制終了ガードを抑止する
\term{契約ホールド}に置く。ホールドはステップ上限の 10 ステップ前
に失効し、ドゥームループ防止のバックストップを保ち、警告のない出
力がこれを解除する。

\subsubsection*{強制終了ガード}

ループは軌跡の両端を守る。早すぎる成功宣言に対しては、ノードは
\texttt{MIN\_TOOL\_CALLS} 回のツール呼び出し前には終了できず、—
そのツールセットに実行ツールが含まれるときは常に — 何かを実行して
その出力を読む前には終了できない。ReAct エージェントで記録された反
復ループの失敗モード(元の研究は推論失敗の大きな割合を、再発行され
たほぼ同一の行動に帰している \cite{yao2023react})に対しては、ルー
プが \term{強制終了}する:ノードが出力を読まれた完了ジョブを持ち、
少なくとも 5 ステップが経過し、契約ホールドが有効でなくなれば — あ
るいは最小ガードを満たした上で上限直前の最後の 3 ステップでは — ツ
ール呼び出し要件が外され、モデルの次の出力はその最終構造化レポート
でなければならなくなる。反復はツールレベルでも減衰される(アイディ
ア生成ツールは最初の成功呼び出しの後に抑止される)し、注入された
ユーザメッセージはメッセージブロック境界へ遅延されるため、再生され
る会話がツール呼び出しをその応答から分割することはない。

\subsection{系譜と再現性}\label{sec:exploration:lineage}

系譜とは最良ノードの先祖の連鎖であり、\texttt{tree.json} の
\texttt{lineage} キーとしてチェックポイントへ書かれる。
\texttt{LineageState} データクラスは BFTS が継続判定に用いる移
動統計を保持し、跨チェックポイントの歩行器
\texttt{walk\_ancestor\_ckpts} は実行を跨ぐ親ポインタを提供す
る。アイディアプールは \texttt{get\_idea\_pool\_for\_ckpt} 経由
で継承され、ランキング付きルートアイディアの選択は
\texttt{RootChoice} データクラスに記録されて事後監査可能となる
(プロンプト: \cref{prompt:root_idea_selector})。

再現性は三つの不変量にかかる。第一に、すべての乱択操作はノード
識別子から自身を播種する。第二に、すべてのツール呼び出しは引数
と出力を該当ノードの作業ツリーへ書き、決して親へは書かない —
これは MCP 層の \texttt{cow\_node\_id} チェックで強制される。第
三に、コスト台帳が \texttt{CostTracker.init} でノードあたりのド
ル額を付与するため、同一プロンプトでの予算検査は決定的となる。

\paragraph{既定はオフライン。}
再現性は探索ループが何を読めるかも規定する。既定では各ノードの
エージェントループは\emph{オフライン}で走る — Web アクセススキ
ルは論文執筆と再現のフェーズにゲートされており、ノードは探索の
最中に時々刻々変わるライブ検索結果を参照できないため、ノードご
との軌跡はその日の Web の状態に依存しない。文献調査は依然行われ
るが、より早く枠の外で — 探索開始前、構想時に行う有界な調査とし
て行われる(\cref{sec:exploration:ideation})。ライブ情報が真に必
要な実行は、オプトイン
(\code{bfts.allow\_web})で探索中のノードエージェントへ Web ア
クセスを開放できる;その場合、実行はチェックポイントに非再現軌
跡マーカーを書き込むため、後の再現が暗黙に厳密扱いされることは
ない。いずれにせよ決定性の契約は誠実に保たれる:既定の実行はバ
イト単位で再現可能であり、再現性をライブ情報と引き換えにする実
行は、その取引を自身の来歴に記録する。

\subsection{計算された証拠のための系譜連鎖}\label{sec:exploration:lineage-chaining}

実行が支えようとする主張のすべてが、新規に測定されるわけではな
い。一部の証拠は、すでに存在する測定から\emph{計算}される —
先行する probe 群に対するパラメータフィット、そのフィットの
held-out 検証、構成間のモデルに基づく選択である。こうした主張は、
互いに独立な木ノードでは構造的に到達不能であることが判明した:
実際の実行では、フィッティングや検証に向けられた子であっても、
測定を持たない親から拡張されると単発 probe の再実行へ退行した。
「すでに手元にあるデータから計算する」が可視の選択肢として一度
も現れない — 子はデータを持たず、どのノードが持つかをセレクタに
教えるものが何もない — ためである。

この間隙を塞ぐ機構は、木が既に持つ親$\to$子の作業ディレクトリ継
承(子は親の作業ツリーの複製内で実行される;
\cref{sec:paper:evidence} も同じ事実に依拠する)のみを、拡張の
両側に 1 つずつのヒントとして用いる。親側では、拡張対象選択の目
標に付加される主張カバレッジブロック
(\cref{sec:exploration:expansion})が、契約証拠名を作業ツリーに
最も多く保持するノードを名指す系譜ヒントを得るため、フィットや検
証を担う子は、その入力を既に保持するノードから優先的に拡張され
る。子側では、継承したディレクトリに系譜上の測定ファイルが含まれ
る子に対し、開始時(\cref{sec:exploration:react} の継承データノー
ト)に、どのファイルが存在しどの正確な契約名を含むかが伝えられ、
基礎にある実験を再実行するのではなく、それらのファイルから導出証
拠を\emph{計算する}よう指示される。

ノード境界を越えるのは名前だけである:ファイル名と測定名であり、
値も、兄弟ノードの結論も、決して越えない。これらのヒントは契約
自体と同種の実行レベルの調整メタデータであるため、
\cref{sec:exploration:memory} のブランチ隔離は保たれ — 欠陥のあ
るブランチの数値が兄弟の推論へ漏れることは依然としてできない —
その一方で、計算された証拠の連鎖(フィット、次いで検証、次いで
選択)は、繰り返しの probe へ潰れる代わりに単一の系譜に沿って実
行できる。

\subsection{プラットフォームプローブ}\label{sec:exploration:platform-probe}

研究計画は、その科学について誤る前に、そのプラットフォームについ
て誤りうる。ある実際の実行では宣言済み主張がハードウェアカウンタ
のプロファイルを要求したが、想定されたプロファイラは実験が走った
計算区画から検証可能なほど明らかに不在だった:それらの主張は恒久
的に充足不能であり、最終ゲートは — 正しく — どのノードも決して生
み出せない証拠ゆえに論文をブロックした。したがってプラットフォー
ムの実現可能性は、想定されるのではなく\emph{測定}されねばならない
— しかも、下流の何かが測定語彙にコミットする前に、一度測定されね
ばならない。

実行がバッチ区画とともに設定されているとき、実行開始は一度きりの
\term{プラットフォームプローブ}を発行する:計算区画自体で実行され
る短いスケジューラジョブで、小さな設定可能な測定ツール一覧(プロ
ファイラ、ハードウェアカウンタユーティリティ、その類)の利用可能
性を名前で確認し、マシンアーキテクチャを記録する。結果はチェック
ポイント生成物(\texttt{platform\_capabilities.json})としてキャッ
シュされ、実行内で再プローブされない。このプローブは設計上ベスト
エフォートである:区画未設定、スケジューラ欠落、あるいはタイムアウ
トを超えるキュー待ちは、いずれも何も書かない記録済みスキップへ降
格する — プローブ失敗はすべての消費者を無制約のままにし、下流の何
も変えない。

プローブされた事実はそのうえで — データとして、決して組み込みのハ
ードウェア知識としてではなく — プラットフォーム盲目な計画が生じか
ねない 3 つの層で中継される。\emph{アイディア}時にはそれらが研究
トピックに織り込まれ、構想段階(\cref{sec:exploration:ideation})は
プラットフォームが取れる測定だけを計画する:実現可能性は下流で繕う
のではなく源で固定される。\emph{契約}時には、選ばれたアイディアの
プランからメトリック契約を導く主張抽出器がそれらを検証済みの能力
ノートとして受け取り、契約が不在のツールに依存する証拠を決して宣言
しないようにする — 契約は一度鋳造されると凍結される
(\cref{sec:paper:contract})ため、二重に重要である。\emph{ノード}時
には、pin された義務メッセージ(\cref{sec:exploration:react})が検証
可能なほど不在のツールを列挙するプラットフォームノートを運び、実行
中のエージェントが欠けたコマンドを苦労して発見するために ReAct ステ
ップを費やさないようにする。\emph{どの}ツールがプローブに値するかの
知識は HPC スキルの設定可能な一覧に存在し、コアのオーケストレータと
ゲートはドメインに依存しないまま、測定されたものを中継するだけであ
る。

\subsection{研究メモリ}\label{sec:exploration:memory}

ノードは孤立して走るわけではない:各ノードは祖先が確立したもの
を継承し、その代わりに子孫と論文執筆器のために記録を残す。ARI
はこの記録を実行ごとの\term{研究メモリ}に保持する。これは自由形
式のログでは満たせない二つの要件に従う。第一に\emph{ブランチ隔
離}:ノードは自身の祖先のメモリのみを読め、兄弟や無関係なチェッ
クポイントのものは決して読まないため、二つの並行する手がかりは
互いを汚染できない — \cref{sec:exploration:lineage} の系譜を定
める祖先スコープ述語と同一である。第二に\emph{基盤づけ可能性}:
メモリは最終的に論文を支えるため、エントリは結果を単に言い直す
のではなく結果の背後にある証拠を指し示せねばならない。そうして
主張はそれを支える生成物まで遡れる。

\subsubsection*{証拠への索引としてのメモリ}

ノードが何を生成したかの唯一の真実の源は、その構造化自己レポー
ト — ノード完了時に書かれる \texttt{node\_report.json} である
(\cref{sec:paper:evidence} がそのフィールドを詳述する)。メモリ
エントリはこれらのフィールドを複製しない。短い検索可能テキスト、
\term{種別}、元のノードレポートへのポインタ、そして結果が依拠す
る特定の生成物 — ログ、ソースファイル、出力ファイル — へのコンテ
ンツハッシュ参照を持つ。したがってメモリエントリは証拠そのもの
ではなく\emph{証拠への索引}である:引用する生成物はいつでも再ハ
ッシュして、なお存在し一致することを確認でき、各参照を検証済み・
欠落・不一致に分類する整合性検査となる。\term{種別}は固定語彙 —
観測、実験結果、失敗事例、手順、内省、生成物要約、論文主張、再
現性イベント — から取られ、読み手は必要な記録を正確に要求できる:
あるブランチ沿いの失敗のみ、あるいは生成物に裏づけられた結果の
み、というように。

\subsubsection*{ノード終端での統合}

ノードが完了すると、決定的な\term{統合}ステップがそのノードレポ
ートを読み、型付きエントリを発する:測定を生んだノードには、測
定した生成物への来歴参照を伴う実験結果エントリを;失敗したノー
ドには失敗事例エントリを。このステップは言語モデルを呼ばないた
め、同じノードレポートは常に同じエントリへ統合される — 統合は記
録された状態の純粋関数であり、チェックポイントの再実行はそれが
生んだメモリを再現する。これはエージェントの動作ではなくループ
フックである:探索エージェントが「メモリに保存せよ」と指示され
ることはなく、実際エージェントが自発的に想起ツールへ手を伸ばす
ことはない。ループが依存するメモリ読み出しはループ自身が決定的
に行う — これが、非決定的な検索をクリティカルパスに置かずにメモ
リへ来歴を担わせる仕掛けである(\cref{sec:discussion:p2})。

\begin{figure}[t]
  \centering
  \resizebox{0.92\columnwidth}{!}{% F18: verifiable research memory as a TYPED INDEX OVER EVIDENCE (not a flow).
% The store holds typed records that point, by content hash, at the artefacts
% that ground them; reproduction outcomes are appended and folded to a latest
% status; an admissibility predicate usable(mu) gates which records may become
% paper claims, so the SAME store yields two projections: a working context for
% exploration and a verified context for the paper.
% style.tikzstyles is loaded by shared/preamble.tex / the standalone wrapper.
\begin{tikzpicture}[node distance=1.5cm and 2.2cm]

  % --- source of truth feeding the store -----------------------------
  \node[artifact wide]                          (report) {\figlabel{F18.report}};
  \node[block emph big, below=1.3cm of report]  (store)  {\figlabel{F18.store}};

  % --- what the records index, and what they accrete -----------------
  \node[block native big, left=of store]        (artifacts) {\figlabel{F18.artifacts}};
  \node[block native big, right=of store]       (repro)     {\figlabel{F18.repro}};
  % record vocabulary: a legend in the empty space under the artefacts
  \node[block aux, below=1.3cm of artifacts]    (kinds)     {\figlabel{F18.kinds}};

  % --- two projections of the SAME store -----------------------------
  \node[block native big, below=2.4cm of store]              (workctx) {\figlabel{F18.workctx}};
  \node[guard, below right=1.4cm and 1.7cm of store]         (usable)  {\figlabel{F18.usable}};
  \node[block emph big, below=1.1cm of usable]               (vctx)    {\figlabel{F18.vctx}};

  % --- edges ----------------------------------------------------------
  \draw[arrow]        (report) -- node[edge label] {\figlabel{F18.consolidate}} (store);
  \draw[arrow dashed] (store)  -- node[edge label] {\figlabel{F18.refs}}        (artifacts);
  \draw[arrow]        (repro)  -- node[edge label] {\figlabel{F18.fold}}        (store);
  \draw[arrow dashed] (store)  -- (kinds);
  \draw[arrow]        (store)  -- node[edge label] {\figlabel{F18.forexpl}} (workctx);
  \draw[arrow]        (store)  -- (usable);
  \draw[arrow]        (usable) -- node[edge label] {\figlabel{F18.forpaper}} (vctx);

\end{tikzpicture}
}
  \caption{検証可能な研究メモリ。各ノードの真実の源である
           \texttt{node\_report.json} の決定的統合が、生成物(ログ・
           出力・ソース)を内容ハッシュで\emph{インデックス}する
           \emph{型付き}格納庫を満たす。再現イベントは追記され最新
           状態へ畳み込まれ、許容述語 $\mathrm{usable}(\mu)$
           (\cref{eq:usable})がどのレコードを検証済み文脈へ通すか
           を制御する。同じ格納庫が二つの読み手を養う:次ノード開始
           時に注入される\emph{作業文脈}(探索の継承経路)と、論文の
           定量的主張を基盤づける\emph{検証済み文脈}
           (\cref{sec:paper:evidence})である。基盤のない内省は探索
           を導くことはあっても論文本文を導くことはない。}
  \label{fig:research-memory}
\end{figure}

\subsubsection*{二つの読み手:作業文脈と検証済み文脈}

格納庫は二つの消費者を養い(\cref{fig:research-memory})、両者の
区別が設計の核心である。各ノードの\emph{開始}時、ループは\term{作
業文脈}を注入する:実験の固定核 — 目標、対象メトリック、ハード
ウェア — と、祖先が記録した結論である。よって子は集約テキストダ
ンプからではなく、系譜がすでに確立したものから、決定的かつ自身
のブランチに限定して始まる。\emph{論文}時には、パイプラインが最
良ノードの系譜から\term{検証済み文脈}を構築する:各結果の再現性
イベントを最新状態へ畳み込み、再現された結果が単に生成物に基盤づ
けられた結果より上位に、後者がさらに基盤のない内省より上位になる
ようエントリを順位づける。エントリ $\mu$ が定量的な論文主張とし
て\emph{許容}されるのは、それが生成物に基盤づけられ、かつ再現の
試みに反証されていないときに限る:
\begin{equation}\label{eq:usable}
  \mathrm{usable}(\mu) \;=\;
    \mathrm{grounded}(\mu)\;\wedge\;
    \bigl(\mathrm{repro}(\mu)\neq\mathrm{failed}\bigr).
\end{equation}
基盤のない内省は作業文脈を通じて探索を導けるが、\cref{eq:usable}
がそれを論文本文から締め出す:執筆器の定量的主張は許容され再現を
優先したエントリのみに依拠する。これは証拠アセンブラ
(\cref{sec:paper:evidence})のメモリ側の対応物である;アセンブラ
が執筆器の読める生成物\emph{バイト}を決めるのに対し、検証済み文
脈はどの\emph{結果}が主張になれるかを決める。

\subsubsection*{追記イベントとしての再現性}

過去のエントリはバイト安定 — ノードは自身の識別子のもとでのみ書
き、既存エントリは決して上書きされない — であるため、結果の再現
性状態は後で再実行されても上書きされない。代わりに再実行は再現性
イベントエントリを\emph{追記}し、読み手は対象ごとのイベントを最
新状態へ畳み込む。再現に失敗した結果は、それを有望と記録した歴史
を書き換えることなく \cref{eq:usable} の主張集合から降格され —
そして正直に限界として報告できる。

\subsection{構想:ルートの播種}\label{sec:exploration:ideation}

木のルートノードはコードから始まるわけではない;何を調査するかを
決めることから始まる。ルートのエージェントはアイディア段階を走ら
せる:有界な文献調査に続いて多エージェントのアイディア生成を行い、
選択されたアイディア(チェックポイント内の \texttt{idea.json})へ
収束する。そのプランが実行のメトリック契約
(\cref{sec:paper:contract})を生む。調査はここで、構想時に行われ
るため、探索自体はオフラインに保てる(\cref{sec:exploration:lineage}):
調査は Semantic~Scholar からライブの文献を取得し — キーワード検索
を引用グラフの 2 ホップ走査で豊かにする — 続く討議は取得した抄録
の上で働く。

構想は意図的に多エージェントである。単一エージェントの、アーカイ
ブに条件づけられたブレインストーミングは、狭いアイディア分布へ崩
落することが知られている — AI~Scientist はそのアイディア生成が
「異なる実行間で、さらにはモデル間でさえ、しばしば非常に似たアイ
ディアをもたらす」と報告する \cite{lu2024ai} — 一方で VirSci
\cite{su2024many} は科学者チーム(共同研究者選択、円卓討論、アイ
ディア生成、新規性投票)をシミュレートし、多エージェントシステム
が単一エージェントのエグゼクティブを、平均で、現代研究との整合性
で $+13.8\%$、潜在的影響で $+44.1\%$ 改善すると報告する。

ARI は VirSci を二つの忠実度でラップする。デフォルト経路はアイディ
アスキル内に VirSci の討論フローを再実装し、利用可能な場合は
vendored の VirSci プロンプトテンプレートを直接用い、検索をライブ
調査に再ベース化し、モデル呼び出しを ARI の LLM 層経由でルーティン
グする;候補アイディアは VirSci 流の自己評価(新規性、実現可能性、
明瞭性)を伴い、新規性で重み付けされた和でランキングされる。オプ
トインの \term{VirSci-live} モードは代わりに、vendored の VirSci
エンジン自体を未改変で駆動する(\cref{fig:virsci}):Semantic-Scholar
由来の凍結スナップショット — SPECTER2 検索インデックスと鮮度でラン
キングされた共著グラフを伴う論文コーパス — がチェックポイント下に
実体化され、VirSci 自身の \texttt{select\_coauthors} がそのグラフ上
に科学者チームを形成し、その $K$ ラウンドの \texttt{generate\_idea}
討議が候補を生む。こうして ARI が構築の基盤とする多エージェント機構
は、言い換えられたものではなく、現在の文献に対して実行される、公刊
されたそのものである;スナップショットの凍結が討議基盤を再現可能に
保ち、ライブ経路はいかなる失敗時にも再実装ループへ降格するため、こ
れを有効化しても実行が退行することは決してない。

二つの実行レベルの事実が、両経路で討議に織り込まれる。プローブ
(\cref{sec:exploration:platform-probe})が検証したプラットフォーム
制約が研究トピックに付加され、トピックを消費するすべてのプロンプ
トがプラットフォームの取れる測定だけを計画する;そして派生実行では
先祖チェックポイントのアイディアプール(\cref{sec:exploration:lineage})
が読み取り専用で注入され、チームは系譜が既に提案したものを再導出す
る代わりに、それを洗練・拡張、あるいは明示的にピボットする。候補間
のランキング付き選択は事後監査のため記録され(\cref{sec:exploration:lineage})、
実行の出発点をその後のどの拡張とも同じく追跡可能に保つ。
